% ============================================================
% Formulario, argomento 12: PDF condizionate e funzioni di variabili continue (dispensa, capitolo 12)
% ============================================================
\argomento{PDF condizionate e funzioni di variabili continue}{Dispensa: capitolo 12}[PDF condizionate e
  funzioni]

\voce{pdfcond}{PDF e CDF condizionate a un evento}{$P(A) > 0$}
  {F_{X \mid A}(x) = P(X \leq x \mid A) = \frac{P(\{X \leq x\} \cap A)}{P(A)} \qquad
   f_{X \mid A}(x) = \frac{\dd}{\dd x} F_{X \mid A}(x)}
  {\variante{$A = \{a \leq X \leq b\}$}{f_{X \mid A}(x) = \frac{f_X(x)}{F_X(b) - F_X(a)} \ \text{ per } a \leq x \leq b,
                                       \quad 0 \ \text{altrove}}}

\voce{totalepdf}{Probabilità totale per le densità}{$\{E_m\}$ partizione di $\Omega$}
  {f_X(x) = \sum_{m=1}^{M} f_{X \mid E_m}(x)\, P(E_m) \qquad F_X(x) = \sum_{m=1}^{M} F_{X \mid E_m}(x)\, P(E_m)}
  {\variante{e per le medie}{E[X] = \sum_{m=1}^{M} E[X \mid E_m]\, P(E_m), \quad
                             E[X \mid E_m] = \int_{-\infty}^{\infty} x\, f_{X \mid E_m}(x) \dd x}}

\voce{monotona}{Densità di $Y = g(X)$ con $g$ monotona}{$x(y) = g^{-1}(y)$, \ $\dot g$ derivata di $g$}
  {f_Y(y) = \frac{f_X\big(x(y)\big)}{\abs{\dot g\big(x(y)\big)}}}
  {\variante{trasformazione lineare, $Y = aX + b$}{f_Y(y) = \frac{1}{\abs{a}}\, f_X\left( \frac{y - b}{a} \right)}
   \variante{la probabilità si conserva}{f_Y(y) \dd y = f_X(x) \dd x}}

\voce{nonmonotona}{Densità di $Y = g(X)$ con $g$ non invertibile}{$x_1(y), \ldots, x_M(y)$ soluzioni di $g(x) = y$}
  {f_Y(y) = \sum_{i=1}^{M(y)} \frac{f_X\big(x_i(y)\big)}{\abs{\dot g\big(x_i(y)\big)}}}
  {\variante{il quadrato, $Y = X^2$}{f_Y(y) = \frac{f_X(\sqrt{y}) + f_X(-\sqrt{y})}{2\sqrt{y}}\, u(y)}}

\voce{cdfuniforme}{La CDF rende uniforme}{$F_X$ strettamente crescente}
  {Y = F_X(X) \sim \mathcal{U}(0, 1)}
  {}

\voce{inversione}{Metodo dell'inversione}{simulare qualsiasi densità}
  {U \sim \mathcal{U}(0, 1) \;\Rightarrow\; Z = F_X^{-1}(U) \ \text{ha CDF } F_X}
  {\variante{esponenziale}{Z = -\frac{1}{\lambda} \ln(1 - U) \sim \mathcal{E}(\lambda)}}

\voce{arcoseno}{Sinusoide con fase a caso}{$X \sim \mathcal{U}\left(-\frac{1}{2}, \frac{1}{2}\right)$, \ $A > 0$}
  {Y = A \cos(2\pi X + \varphi): \quad f_Y(y) = \frac{1}{\pi \sqrt{A^2 - y^2}}, \quad y \in (-A, A)}
  {\variante{è la distribuzione dell'arcoseno, a forma di «U»}{}}

\voce{conversionead}{Conversione A/D a $R$ bit}{$M = 2^R$ livelli, soglie $x_1 < \cdots < x_{M+1}$}
  {Y = y_i \ \text{ se } x_i \leq X < x_{i+1} \qquad p_Y(y_i) = F_X(x_{i+1}) - F_X(x_i)}
  {}
